\chapter{Contexte et revue de littérature}
\label{chapitre:contexte}

% =====================================================================
\section{Cadre général}
% =====================================================================

La boiterie bovine est un problème de santé majeur en élevage laitier, avec des conséquences
directes sur le bien-être animal, la productivité et la rentabilité des exploitations
\citep{rodrigues2017lameness}.
Des études épidémiologiques rapportent des prévalences allant de 20\% à plus de 50\% selon
les régions, les systèmes d'élevage et les critères de diagnostic utilisés
\citep{whay2003assessment,cook2003prevalence}. Les pathologies du pied, qu'elles soient
infectieuses (\textit{dermatite digitale}, \textit{phlegmon interdigité}) ou non
infectieuses (\textit{ulcère de la sole}, \textit{maladie de la ligne blanche}),
comptent parmi les principales causes de boiterie chez les vaches laitières
\citep{merck_hoof_origin_lameness_cattle,flower2006effect}.

\subsection{Impact économique}
L'impact économique de la boiterie est substantiel et multidimensionnel. Le
Tableau~\ref{tab:impact_economique} synthétise les principales composantes du coût
\textit{par cas clinique}, c'est-à-dire par épisode de boiterie diagnostiqué chez une
vache au cours d'une lactation, d'après les estimations publiées dans la littérature.

\begin{table}[H]
    \centering
    \caption{Composantes du coût économique de la boiterie par cas clinique.}
    \label{tab:impact_economique}
    \begin{tabular}{lcc}
        \toprule
        \textbf{Composante} & \textbf{Estimation} & \textbf{Source} \\
        \midrule
        Perte de production laitière & 270--574 kg/lactation & \cite{green2002financial} \\
        Coût vétérinaire direct & 25--90 USD/cas & \cite{bruijnis2010assessing} \\
        Dégradation de la fertilité & +30--50 jours ouverts & \cite{green2002financial} \\
        Réforme anticipée & 75--300 USD/cas & \cite{bruijnis2010assessing} \\
        \midrule
        \textbf{Coût total estimé} & \textbf{75--300+ USD/cas} & \\
        \bottomrule
    \end{tabular}
\end{table}
\vspace{-0.9\baselineskip}

Quatre postes principaux composent le coût d'un cas clinique: la perte de production
laitière (270 à 574~kg par lactation), liée à la baisse de l'ingestion et de l'activité;
le coût vétérinaire direct (25 à 90~USD par cas), qui inclut la consultation, le parage,
les traitements et le matériel de soin; la dégradation de la fertilité, mesurée par une
hausse de 30 à 50 \textit{jours ouverts}, soit l'intervalle entre le vêlage et la
conception suivante, avec un décalage de la lactation suivante; et la réforme anticipée
(75 à 300~USD par cas), liée à la perte de valeur résiduelle d'une vache écartée avant la
fin de sa carrière productive. Ces postes se cumulent selon la sévérité et la chronicité
de l'épisode. À l'échelle d'un troupeau de 100~vaches avec une prévalence de 25\,\%, le
coût annuel cumulé est estimé entre 1\,900 et 7\,500~USD pour les seuls cas incidents.
Au-delà de l'impact financier, la boiterie soulève un enjeu éthique de bien-être animal~:
la douleur associée est reconnue comme l'une des principales sources de souffrance
chronique chez les vaches laitières \citep{whay2003assessment,broom2017animal}. Ces conséquences
justifient une compréhension fine des mécanismes de boiterie, afin d'orienter la
prévention, la détection et la prise en charge.

\subsection{Étiologie et facteurs de risque}
Les causes de la boiterie sont multifactorielles. La Figure~\ref{fig:etiologie} les
organise en quatre catégories~: causes directes infectieuses ou non infectieuses, et
facteurs modulateurs environnementaux ou individuels. Les causes directes regroupent des
affections bactériennes, comme la \textit{dermatite digitale} ou le \textit{phlegmon
interdigité}, et des lésions mécaniques, comme l'\textit{ulcère de la sole} ou la
\textit{maladie de la ligne blanche}. Les facteurs environnementaux (sol, hygiène,
humidité) modulent l'exposition, tandis que les facteurs individuels (parité, âge, stade
de lactation) modulent la susceptibilité. Cette structure multifactorielle explique
qu'aucun indicateur unique ne suffise et justifie une approche fondée sur plusieurs signaux
pour la détection précoce
\citep{merck_hoof_origin_lameness_cattle,dairynz_preventing_managing_lameness_guide_2025}.

\begin{figure}[H]
    \centering
    \resizebox{0.85\textwidth}{!}{%
    \begin{tikzpicture}[
        cat/.style={rectangle, draw=black!70, rounded corners=3pt, minimum height=0.9cm,
                    minimum width=3.2cm, align=center, font=\small, inner sep=10pt},
        factor/.style={rectangle, draw=black!40, rounded corners=2pt, minimum height=0.8cm,
                      minimum width=2.7cm, align=center, font=\small, inner sep=6pt},
        arrow/.style={-{Stealth[length=2mm]}, thick, black!60},
    ]
        % Titre central
        \node[cat, fill=rawred!15, minimum width=4cm, font=\small\bfseries] (boiterie)
            {Boiterie bovine};

        % Catégories
        \node[cat, fill=ifblue!12, above left=1.5cm and 1cm of boiterie] (infect)
            {Pathologies\\[-7pt]infectieuses};
        \node[cat, fill=loforange!12, above right=1.5cm and 1cm of boiterie] (noninfect)
            {Pathologies\\[-7pt]non infectieuses};
        \node[cat, fill=rulegreen!12, below left=1.5cm and 1cm of boiterie] (enviro)
            {Facteurs\\[-7pt]environnementaux};
        \node[cat, fill=profgray!12, below right=1.5cm and 1cm of boiterie] (indiv)
            {Facteurs\\[-7pt]individuels};

        % Liens
        \draw[arrow] (infect) -- (boiterie);
        \draw[arrow] (noninfect) -- (boiterie);
        \draw[arrow] (enviro) -- (boiterie);
        \draw[arrow] (indiv) -- (boiterie);

        % Sous-facteurs infectieux
        \node[factor, fill=ifblue!6, above=0.9cm of infect, xshift=-2.0cm] (dd)
            {Dermatite digitale};
        \node[factor, fill=ifblue!6, above=0.9cm of infect, xshift=2.0cm] (phleg)
            {Phlegmon interdigité};
        \draw[black!30] (dd) -- (infect);
        \draw[black!30] (phleg) -- (infect);

        % Sous-facteurs non infectieux
        \node[factor, fill=loforange!6, above=0.9cm of noninfect, xshift=-2.0cm] (ulcere)
            {Ulcère de la sole};
        \node[factor, fill=loforange!6, above=0.9cm of noninfect, xshift=2.0cm] (ligne)
            {Maladie ligne blanche};
        \draw[black!30] (ulcere) -- (noninfect);
        \draw[black!30] (ligne) -- (noninfect);

        % Sous-facteurs environnementaux
        \node[factor, fill=rulegreen!6, below=0.9cm of enviro, xshift=-2.0cm] (sol)
            {Type de sol};
        \node[factor, fill=rulegreen!6, below=0.9cm of enviro, xshift=2.0cm] (hygiene)
            {Hygiène, humidité};
        \draw[black!30] (sol) -- (enviro);
        \draw[black!30] (hygiene) -- (enviro);

        % Sous-facteurs individuels
        \node[factor, fill=profgray!6, below=0.9cm of indiv, xshift=-2.0cm] (parite)
            {Parité, âge};
        \node[factor, fill=profgray!6, below=0.9cm of indiv, xshift=2.0cm] (lactation)
            {Stade de lactation};
        \draw[black!30] (parite) -- (indiv);
        \draw[black!30] (lactation) -- (indiv);
    \end{tikzpicture}%
    }
    \caption{Étiologie multifactorielle de la boiterie bovine.}
    \label{fig:etiologie}
\end{figure}

\subsection{Importance de la détection précoce}
La littérature sur la boiterie bovine souligne l'importance d'une détection précoce pour
limiter la dégradation du bien-être animal et les impacts de production. Les ressources
vétérinaires de référence et les guides techniques convergent sur deux idées:
une surveillance régulière est indispensable, et la prise en charge doit être
structurée et rapide selon la sévérité du cas
\citep{merck_lameness_overview,dairynz_treating_lameness}.

La détection précoce (score locomoteur 2 sur une échelle de 1 à 5) ouvre la voie au
traitement des cas avant qu'ils ne deviennent chroniques. Les études montrent qu'un
retard de traitement de quelques jours peut multiplier la durée de récupération et le
coût associé \citep{green2002financial}. Les capteurs portés et les systèmes de mesure de
la démarche permettent de quantifier objectivement des modifications d'activité ou de
locomotion associées à la boiterie
\citep{chapinal2010measurement,pastell2009measurement}. Le délai avec lequel ces
modifications deviennent détectables avant une observation clinique dépend toutefois du
capteur, de la définition du cas et du protocole; il ne peut donc pas être généralisé à
partir de ces travaux. En pratique, les boiteries légères restent difficiles à repérer de
façon régulière par la seule observation visuelle.

Des revues de synthèse en élevage de précision montrent aussi que la boiterie fait partie
des cas d'usage les plus importants, mais également des plus difficiles à stabiliser, car
les signaux disponibles restent indirects, multifactoriels et fortement dépendants du
contexte d'élevage \citep{palma2025scoping,michelena2025plf_review}.

% =====================================================================
\section{Approches conventionnelles~: observation et notation locomotrice}
% =====================================================================

Les approches conventionnelles reposent principalement sur l'observation humaine et les
grilles de notation locomotrice. La grille proposée par \citet{sprecher1997lameness}
reste l'une des plus utilisées en pratique. Elle attribue un score de 1
(\textit{normal}) à 5 (\textit{sévèrement boiteux}) en se basant sur la posture du dos
et la démarche \citep{flower2006gait}. La fiche opérationnelle de l'AHDB propose une
grille complémentaire adaptée au contexte britannique et associe le score~2 à une
intervention rapide, généralement dans les 48~heures
\citep{ahdb_mobility_score_sheet_48h_score2}.

Les approches conventionnelles restent pertinentes en élevage car elles sont simples à
déployer et directement interprétables par les équipes terrain. Elles présentent toutefois
plusieurs limites opérationnelles. La première tient à la subjectivité inter-observateurs: même avec
des grilles standardisées, la concordance entre évaluateurs peut varier significativement
\citep{flower2006effect,march2007interobserver}. La deuxième concerne le temps requis:
dans les grands troupeaux, l'observation systématique de chaque animal à chaque traite est
difficile à maintenir. Enfin, les boiteries légères (score~2) restent difficiles à repérer
de manière reproductible par la seule observation visuelle
\citep{whay2003assessment,march2007interobserver}.

Pour dépasser ces limites, le recours à des signaux capteurs vient compléter la
surveillance humaine, sans la remplacer.

\subsection{Lecture opérationnelle du score locomoteur}
Au-delà de sa fonction descriptive, l'évaluation locomotrice est aussi un outil de
priorisation. Dans la fiche opérationnelle de l'AHDB, le passage du score~1
(\textit{normal}) au score~2 (\textit{boiterie légère}) appelle une intervention rapide,
généralement dans les 48~heures
\citep{ahdb_mobility_score_sheet_48h_score2}. Cette logique explique
pourquoi un système automatisé peut viser une fonction plus prudente que la reproduction
d'une note clinique: signaler une dégradation comportementale suffisamment tôt pour
déclencher une vérification sur le terrain.

La Figure~\ref{fig:sprecher_scale} synthétise la lecture opérationnelle du score
locomoteur: le passage vers le score~2 constitue la zone d'intérêt pour une alerte
précoce, car il marque l'entrée dans une situation où l'observation humaine ou l'examen
du pied devient prioritaire
\citep{sprecher1997lameness,ahdb_mobility_score_sheet_48h_score2}.

\begin{figure}[H]
    \centering
    \begin{tikzpicture}[
        score/.style={rectangle, draw=black!60, rounded corners=3pt, minimum width=2.45cm,
                      minimum height=1.7cm, align=center, inner sep=8pt, font=\small},
        arrow/.style={-{Stealth[length=2.5mm]}, thick, black!55}
    ]
        \node[score, fill=rulegreen!12] (s1) at (0,0.2) {\textbf{Score 1}\\[2pt]Normal};
        \node[score, fill=ifblue!10] (s2) at (3.2,0.2) {\textbf{Score 2}\\[2pt]Léger};
        \node[score, fill=loforange!14] (s3) at (6.4,0.2) {\textbf{Score 3}\\[2pt]Modéré};
        \node[score, fill=rawred!14] (s4) at (9.6,0.2) {\textbf{Score 4}\\[2pt]Marqué};
        \node[score, fill=rawred!25] (s5) at (12.8,0.2) {\textbf{Score 5}\\[2pt]Sévère};

        \draw[arrow] (1.35,0.2) -- (1.95,0.2);
        \draw[arrow] (4.55,0.2) -- (5.15,0.2);
        \draw[arrow] (7.75,0.2) -- (8.35,0.2);
        \draw[arrow] (10.95,0.2) -- (11.55,0.2);

        \draw[decorate, decoration={brace, amplitude=6pt}, thick, ifblue]
            (1.7,1.7) -- (4.7,1.7)
            node[midway, above=8pt, align=center, font=\footnotesize\bfseries]
            {Zone d'intérêt pour\\la détection précoce};
    \end{tikzpicture}

    \caption[Seuil opérationnel du score locomoteur]{Seuil opérationnel du score
    locomoteur pour la détection précoce.}
    \label{fig:sprecher_scale}
\end{figure}

Le Tableau~\ref{tab:score_operationnel} précise la lecture clinique simplifiée et
l'action associée à chaque score, d'après la grille de \citet{sprecher1997lameness} et
les recommandations AHDB \citep{ahdb_mobility_scoring_dairy_cows}.

\begin{table}[H]
    \centering
    \caption{Lecture clinique et opérationnelle simplifiée du score locomoteur.}
    \label{tab:score_operationnel}
    \footnotesize
    \renewcommand{\arraystretch}{1.2}
    \begin{tabular}{>{\raggedright\arraybackslash}p{1.4cm}
                    >{\raggedright\arraybackslash}p{6.0cm}
                    >{\raggedright\arraybackslash}p{5.2cm}}
        \toprule
        \textbf{Score} & \textbf{Lecture clinique simplifiée} & \textbf{Lecture opérationnelle} \\
        \midrule
        1 & Dos plat en station et en marche. & Surveillance usuelle. \\
        2 & Dos plat à l'arrêt, légèrement arqué en marche. & Intervention rapide: zone clé pour la détection précoce. \\
        3 & Dos arqué en permanence. & Vérification prioritaire. \\
        4 & Atteinte locomotrice visible. & Prise en charge urgente. \\
        5 & Appui très limité. & Prise en charge immédiate. \\
        \bottomrule
    \end{tabular}
\end{table}

La lecture par seuil transforme la notation locomotrice en problème d'alerte: il ne
s'agit plus seulement de décrire une boiterie, mais de repérer l'entrée dans une zone
d'intervention. Ce déplacement justifie le recours à des signaux capteurs et à des
méthodes d'apprentissage capables de suivre les changements comportementaux dans le
temps.

% =====================================================================
\section{Approches basées sur des capteurs et apprentissage}
% =====================================================================

Les approches basées sur des capteurs visent à détecter des déviations comportementales
(activité, transitions, dynamique de mouvement) qui peuvent être associées à la boiterie,
sans être spécifiques à cette affection.
La littérature distingue trois familles principales: l'analyse de la démarche par vision
par ordinateur \citep{vanhertem2014evaluation,lavrova2023lameness}, la classification de
comportements par accélérométrie \citep{oleary2020accelerometers,alsaaod2019automatic,
balasso2023behaviour,benaissa2019calving} et la détection d'anomalies sur des indicateurs dérivés
\citep{qiao2021review,chandola2009anomaly}. Les revues
récentes confirment la prédominance des approches supervisées, tout en soulignant le
manque d'études non supervisées en conditions réelles d'élevage
\citep{qiao2021review,oleary2020accelerometers}. Une revue récente de la détection
automatisée par intelligence artificielle confirme cette dynamique tout en pointant des
freins persistants: hétérogénéité des protocoles, dépendance aux étiquettes et validation
encore limitée en conditions réelles \citep{siachos2024automated}. Les approches récentes par
vision et apprentissage profond atteignent de bonnes performances sur vidéo
\citep{myint2024lameness}, mais supposent une infrastructure caméra et des données étiquetées,
deux conditions absentes du corpus accéléromètre non étiqueté utilisé ici. Des jeux de données
publics annotés pour la boiterie commencent certes à apparaître, comme \textit{CowScreeningDB}
\citep{ismail2024cowscreeningdb}, mais ils restent rares, propres à un capteur donné, et ne
couvrent pas le corpus IceQube exploité ici.

La diversité de ces approches peut être organisée selon trois dimensions complémentaires:
\textit{la source de données}, \textit{la stratégie analytique} et \textit{la nature de la
sortie}. Une telle lecture est importante, car elle évite de comparer directement des travaux
qui ne visent pas le même objectif. Un modèle supervisé de classification binaire à partir
d'images et un système d'alerte basé sur des capteurs d'activité n'ont ni les mêmes
contraintes de données, ni le même niveau d'interprétabilité, ni le même usage en ferme.

La Figure~\ref{fig:taxonomie_approches} synthétise cette organisation en trois axes~: la
source de données (de l'observation visuelle aux accéléromètres portés), la stratégie
analytique (règles et seuils, modèles supervisés, détection d'anomalies non supervisée,
approches hybrides) et la sortie opérationnelle (du score locomoteur à l'alerte
nécessitant une vérification). Le présent travail se positionne dans la branche non supervisée et hybride,
encadrée en vert sur la figure.

\begin{figure}[H]
    \centering
    \resizebox{\textwidth}{!}{%
    \begin{tikzpicture}[
        block/.style={rectangle, draw=black!70, rounded corners=3pt, minimum width=4cm,
                      minimum height=1cm, align=center, font=\small\bfseries, inner sep=8pt},
        small/.style={rectangle, draw=black!35, rounded corners=2pt, minimum width=2.8cm,
                      minimum height=0.85cm, text width=2.6cm, align=center, font=\footnotesize, inner sep=6pt},
        arrow/.style={-{Stealth[length=2.5mm]}, thick, black!55}
    ]
        \node[block, fill=ifblue!12] (data) at (0,0) {Sources de données};
        \node[block, fill=loforange!12] (model) at (7,0) {Stratégie analytique};
        \node[block, fill=rulegreen!12] (output) at (14,0) {Sortie opérationnelle};

        \node[small, fill=ifblue!6] (obs) at (-1.8,-2.2) {Observation\\et notation};
        \node[small, fill=ifblue!6] (vision) at (1.8,-2.2) {Vision\\de la démarche};
        \node[small, fill=ifblue!6] (accel) at (-1.8,-4.0) {Accéléromètres\\activité / couchage};
        \node[small, fill=ifblue!6] (multi) at (1.8,-4.0) {Multi-source\\capteurs + contexte};

        \node[small, fill=loforange!6] (rules) at (5.2,-2.2) {Règles\\et seuils};
        \node[small, fill=loforange!6] (sup) at (8.8,-2.2) {Supervisé};
        \node[small, fill=loforange!6] (unsup) at (5.2,-4.0) {Non supervisé\\(anomalies)};
        \node[small, fill=loforange!6] (hyb) at (8.8,-4.0) {Hybride\\score + règles};

        \node[small, fill=rulegreen!6] (score) at (12.2,-2.2) {Score\\locomoteur};
        \node[small, fill=rulegreen!6] (classif) at (15.8,-2.2) {Classe binaire\\boiteux / non};
        \node[small, fill=rulegreen!6] (episode) at (12.2,-4.0) {Épisode\\d'alerte};
        \node[small, fill=rulegreen!6] (alerte) at (15.8,-4.0) {Alerte de\\vérification};

        \draw[arrow] (data) -- (model);
        \draw[arrow] (model) -- (output);

        % Lignes: parent -> rang 1 (directes, courtes)
        \draw[black!30] (data.south) -- ++(0,-0.4) -| (obs.north);
        \draw[black!30] (data.south) -- ++(0,-0.4) -| (vision.north);
        \draw[black!30] (model.south) -- ++(0,-0.4) -| (rules.north);
        \draw[black!30] (model.south) -- ++(0,-0.4) -| (sup.north);
        \draw[black!30] (output.south) -- ++(0,-0.4) -| (score.north);
        \draw[black!30] (output.south) -- ++(0,-0.4) -| (classif.north);

        % Lignes: rang 1 -> rang 2 (verticales simples)
        \draw[black!30] (obs.south) -- (accel.north);
        \draw[black!30] (vision.south) -- (multi.north);
        \draw[black!30] (rules.south) -- (unsup.north);
        \draw[black!30] (sup.south) -- (hyb.north);
        \draw[black!30] (score.south) -- (episode.north);
        \draw[black!30] (classif.south) -- (alerte.north);

        % Encadrement du positionnement de ce mémoire (autour des nœuds unsup et hyb)
        \node[draw=rulegreen, thick, rounded corners=3pt, inner sep=6pt,
              fit=(unsup)(hyb)] (encadre) {};
        \node[font=\footnotesize\bfseries, rulegreen!80!black, below=4pt of encadre]
            {Positionnement de ce mémoire};
    \end{tikzpicture}%
    }
    \caption{Taxonomie des approches de détection automatisée.}
    \label{fig:taxonomie_approches}
\end{figure}

Du point de vue de la stratégie analytique, les méthodes de détection se déclinent
classiquement en deux familles:
\begin{itemize}
    \item[(a)] \textit{Modèles supervisés}: lorsque des labels fiables et suffisamment
    représentatifs sont disponibles. Ces approches peuvent atteindre de bonnes performances
    de classification
    \citep{vanhertem2014evaluation,viazzi2013analysis}, mais leur
    dépendance aux labels limite leur déployabilité dans des contextes où l'étiquetage est
    incomplet ou incertain;
    \item[(b)] \textit{Détection d'anomalies et règles expertes}: lorsque les labels sont
    partiels ou absents. Ces approches ne nécessitent pas d'exemples positifs pour
    l'entraînement et sont donc mieux adaptées aux situations réelles où la boiterie n'est
    pas systématiquement diagnostiquée et enregistrée.
\end{itemize}

Dans le corpus étudié, la supervision clinique exhaustive n'est pas disponible. Cette
contrainte motive une architecture qui combine des mesures de déviation et des règles
temporelles explicites \citep{chandola2009anomaly,palma2025scoping}. L'étude de
\citet{michelena2025comparative} compare plusieurs détecteurs non supervisés sur un autre
usage bovin, la détection des chaleurs; elle soutient l'intérêt d'une comparaison empirique
des algorithmes, mais ne fournit pas de preuve sur la boiterie.

\subsection{Isolation Forest}
L'algorithme Isolation Forest (IF), proposé par Liu et al. \citep{liu2008isolation}, repose
sur le principe que les anomalies sont plus faciles à isoler que les points normaux dans un
espace de variables. Son échantillonnage et ses arbres aléatoires permettent de traiter des
jeux tabulaires sans modèle explicite de la distribution normale. Son implémentation dans
scikit-learn \citep{sklearn_api} fournit ici un comparateur ponctuel reproductible; elle ne
modélise toutefois pas directement la direction ni la durée d'un changement.

La Figure~\ref{fig:if_principle} illustre le principe fondamental d'IF: une anomalie
(point rouge) est isolée en seulement deux coupures aléatoires, tandis que les points
normaux (cluster bleu) nécessitent de nombreuses partitions supplémentaires. Le score
d'anomalie est inversement proportionnel au nombre de coupures.

\begin{figure}[H]
    \centering
    \resizebox{0.65\textwidth}{!}{%
    \begin{tikzpicture}
        % Cadre
        \draw[black!50] (0,0) rectangle (12,6);

        % --- Cluster normal : 15 points groupés à gauche ---
        \foreach \x/\y in {2.0/3.0, 2.3/3.4, 2.6/2.8, 2.2/3.2, 1.9/3.5,
                           2.5/3.0, 2.1/2.7, 2.4/3.3, 2.7/3.5, 1.8/3.1,
                           2.2/2.9, 2.5/3.4, 2.3/2.8, 2.1/3.3, 2.6/3.2} {
            \fill[ifblue] (\x,\y) circle (2.5pt);
        }

        % --- Anomalie : 1 point isolé à droite ---
        \fill[rawred] (9.5,4.5) circle (4pt);

        % --- Coupure 1 : verticale entre cluster et anomalie ---
        \draw[black!70, dashed, thick] (6,0.3) -- (6,5.7);

        % --- Coupure 2 : horizontale, zone droite seulement ---
        \draw[black!70, dashed, thick] (6.3,2.8) -- (11.7,2.8);

        % --- Étiquettes hors figure (marges) ---
        % Cluster
        \node[font=\footnotesize, black] at (2.3,1.5) {Points normaux};
        \node[font=\scriptsize, black!50] at (2.3,1.0) {difficiles à isoler};

        % Anomalie
        \node[font=\footnotesize, black] at (9.5,5.3) {Anomalie};
        \node[font=\scriptsize, black!50] at (9.5,3.7) {isolée en 2 coupures};

        % Numéros des coupures
        \node[font=\scriptsize, black, fill=white, inner sep=2pt] at (6,0.9) {1};
        \node[font=\scriptsize, black, fill=white, inner sep=2pt] at (10.8,2.8) {2};
    \end{tikzpicture}%
    }
    \caption{Principe d'Isolation Forest.}
    \label{fig:if_principle}
\end{figure}

\subsection{Local Outlier Factor}
Le Local Outlier Factor (LOF), proposé par Breunig et al. \citep{breunig2000lof}, estime la
densité locale de chaque point par rapport à celle de ses voisins. Il vise ainsi des points
rares dans leur voisinage, tandis qu'IF repose sur leur facilité d'isolation globale. Le coût
de LOF dépend notamment de la recherche des voisins et de la dimension des données. Dans ce
mémoire, il sert de second comparateur non supervisé; aucune supériorité générale ne lui est
attribuée a priori.

\subsection{Comparaison synthétique IF et LOF}
Le Tableau~\ref{tab:if_vs_lof} synthétise les différences fondamentales entre les deux
détecteurs utilisés dans ce projet.

\begin{table}[H]
    \centering
    \caption{Comparaison des propriétés d'Isolation Forest et de Local Outlier Factor.}
    \label{tab:if_vs_lof}
    \begin{tabular}{>{\raggedright\arraybackslash}p{3.2cm}
                    >{\raggedright\arraybackslash}p{5.0cm}
                    >{\raggedright\arraybackslash}p{5.0cm}}
        \toprule
        \textbf{Propriété}            & \textbf{Isolation Forest}
            & \textbf{Local Outlier Factor} \\
        \midrule
        Principe                      & Isolabilité par partitions
            & Densité locale relative \\
        Type d'anomalie ciblée        & Globale (isolée dans l'espace)
            & Locale (faible densité relative) \\
        Référence fondatrice          & \citet{liu2008isolation}
            & \citet{breunig2000lof} \\
        Dépendance principale         & Nombre d'arbres et contamination
            & Nombre de voisins et contamination \\
        Limite pour ce projet         & Ne modélise pas la persistance
            & Ne modélise pas la persistance \\
        Paramètres critiques          & \texttt{contamination}, \texttt{n\_estimators}
            & \texttt{n\_neighbors}, \texttt{contamination} \\
        \bottomrule
    \end{tabular}
\end{table}

IF fournit un comparateur fondé sur l'isolation globale, tandis que LOF apporte un contraste
fondé sur la densité locale. La sortie principale du mémoire repose sur les détecteurs
temporels décrits au Chapitre~\ref{chapitre:methodologie}.

\subsection{Choix méthodologique et travaux directement pertinents}

La revue présentée ici est narrative et ciblée; elle ne constitue ni une revue systématique
ni une méta-analyse. Les travaux ont été retenus parce qu'ils éclairent directement l'une
des décisions du projet: nature indirecte des signaux comportementaux, absence de labels
cliniques dans le corpus principal, intérêt d'une référence individuelle, nécessité d'une
évaluation temporelle et exigence d'interprétabilité.

Les modèles supervisés et les réseaux profonds peuvent exploiter des motifs complexes lorsque
des signaux suffisamment riches et des labels contemporains sont disponibles
\citep{vanhertem2014evaluation,balasso2023behaviour}. Ces conditions ne sont pas réunies
ici: les variables IceQube sont déjà agrégées par intervalles de 15~minutes et le corpus
principal ne comporte pas de diagnostic locomoteur synchronisé. Une classification clinique
supervisée ne serait donc pas défendable. Les méthodes non supervisées restent utiles comme
points de comparaison, mais elles ne résolvent pas à elles seules la direction ni la
persistance d'un changement. Des travaux récents transposent la détection d'anomalies non
supervisée à des séries capteurs d'élevage bruitées \citep{kakar2024timetector}, ce qui
conforte l'intérêt d'approches sans étiquettes; ils restent toutefois orientés vers la
reconstruction de signaux plutôt que vers une baisse comportementale directionnelle,
persistante et individualisée comme celle que cible HYPO.

Le Tableau~\ref{tab:revue_ciblee} résume les travaux les plus directement reliés aux choix
méthodologiques. Il ne compare pas leurs performances numériques, car les capteurs, les
définitions de cas et les unités statistiques diffèrent.

\begin{table}[H]
    \centering
    \small
    \caption{Travaux directement reliés aux choix méthodologiques.}
    \label{tab:revue_ciblee}
    \begin{tabular}{>{\raggedright\arraybackslash}p{3.0cm}
                    >{\raggedright\arraybackslash}p{4.7cm}
                    >{\raggedright\arraybackslash}p{6.0cm}}
        \toprule
        \textbf{Travail} & \textbf{Constat utile} & \textbf{Conséquence pour ce mémoire} \\
        \midrule
        \citet{alsaaod2019automatic}
        & Les systèmes automatiques reposent sur des modalités et des protocoles hétérogènes.
        & Limiter la sortie à une alerte comportementale et documenter le protocole complet. \\
        \addlinespace
        \citet{oleary2020accelerometers}
        & Les variables accélérométriques associées à la boiterie varient selon le capteur et
          le protocole.
        & Éviter une signature universelle et conserver plusieurs familles de signaux. \\
        \addlinespace
        \citet{lavrova2023lameness}
        & Une étude multi-fermes combine activité, couchage, production et caractéristiques
          individuelles.
        & Comparer chaque vache à sa propre référence et reconnaître l'intérêt futur de
          variables contextuelles. \\
        \addlinespace
        \citet{balasso2023behaviour}
        & Les signaux accélérométriques bruts permettent de modéliser des motifs
          comportementaux complexes.
        & Ne pas transposer ces capacités à des agrégats de 15~minutes qui ne contiennent
          plus la dynamique fine de la démarche. \\
        \addlinespace
        \citet{palma2025scoping}
        & Le déploiement de l'intelligence artificielle en élevage dépend de la qualité des
          données, de l'intégration et de l'usage.
        & Privilégier une sortie traçable, interprétable et destinée à une vérification
          humaine. \\
        \bottomrule
    \end{tabular}
\end{table}

Ces travaux motivent une architecture classique, individualisée et temporelle. IF et LOF
sont conservés comme comparateurs ponctuels de familles différentes; le résultat principal
repose toutefois sur une accumulation directionnelle des baisses et sur la concordance de
plusieurs signaux. L'apprentissage profond n'est pas rejeté en général: il est simplement
hors de portée des données agrégées et non étiquetées utilisées dans ce mémoire.

% =====================================================================
\section{Capteurs et variables comportementales}
\label{sec:technologies_capteurs}
% =====================================================================

Avant d'aborder l'hétérogénéité des protocoles, cette section poursuit un double objectif:
situer d'abord brièvement les grandes familles de capteurs disponibles pour la détection
de boiterie, puis décrire le capteur effectivement retenu dans ce projet (IceQube) ainsi
que les variables comportementales qu'il produit. Les systèmes de surveillance
comportementale reposent sur plusieurs familles de capteurs, dont les caractéristiques
conditionnent directement les variables disponibles et les performances atteignables.

Le Tableau~\ref{tab:technologies_capteurs} présente les principales technologies et leurs
caractéristiques opérationnelles. On distingue trois grandes familles: les capteurs portés
par l'animal, tels que les accéléromètres, les podomètres et les colliers; les capteurs
intégrés à l'infrastructure, comme les tapis de pesée et les systèmes de traite; et les
systèmes sans contact, comme les caméras et la vision par ordinateur. Les accéléromètres
portés constituent la technologie la plus étudiée dans la littérature récente
\citep{oleary2020accelerometers,alsaaod2019automatic}.

\begin{table}[H]
    \centering
    \small
    \caption{Principales familles de capteurs étudiées.}
    \label{tab:technologies_capteurs}
    \begin{tabular}{>{\raggedright\arraybackslash}p{2.4cm}
                    >{\raggedright\arraybackslash}p{2.8cm}
                    >{\raggedright\arraybackslash}p{2.8cm}
                    >{\raggedright\arraybackslash}p{2.2cm}
                    >{\raggedright\arraybackslash}p{2.2cm}}
        \toprule
        \textbf{Technologie} & \textbf{Variables mesurées} & \textbf{Avantages} &
        \textbf{Limites} & \textbf{Études repr.} \\
        \midrule
        Accéléromètre (patte/collier)
        & Activité, pas, transitions couché/debout, durée couchage
        & Continu, individuel, coût modéré
        & Dérive, position variable
        & \cite{alsaaod2012electronic,balasso2023behaviour} \\
        \addlinespace
        Podomètre
        & Nombre de pas, distance parcourue
        & Simple, robuste
        & Signal limité, faible résolution
        & \cite{flower2006effect,mazrier2006reliability} \\
        \addlinespace
        Passerelle sensible à la pression
        & Pressions plantaires et paramètres spatio-temporels de la démarche
        & Objectif, quantitatif
        & Infrastructure fixe, coût élevé
        & \cite{maertens2011gaitwise} \\
        \addlinespace
        Vision par ordinateur
        & Posture, arc du dos, démarche
        & Sans contact, riche en information
        & Sensible aux conditions, calibration
        & \cite{qiao2021review} \\
        \addlinespace
        Système de traite automatisé
        & Production laitière, conductivité, fréquentation
        & Déjà installé, données continues
        & Signal indirect pour la boiterie
        & \cite{lavrova2023lameness,vanhertem2013lameness} \\
        \bottomrule
    \end{tabular}
\end{table}

Dans le cadre de ce mémoire, les données proviennent de capteurs IceQube
(IceRobotics Ltd, Édimbourg, Royaume-Uni), des accéléromètres tri-axiaux conçus
spécifiquement pour le suivi comportemental des vaches laitières. La
Figure~\ref{fig:iceqube_capteur} montre la représentation schématique du capteur IceQube,
dont le boîtier étanche contient l'accéléromètre tri-axial et se fixe à la patte par un
bracelet en plastique muni d'une boucle et de trous d'ajustement. La
Figure~\ref{fig:iceqube_position} illustre son positionnement sur le métatarse de la patte
arrière~; les axes $x$, $y$, $z$ de l'accéléromètre permettent de distinguer la position
debout (patte verticale, $\vec{g}$ aligné avec $z$) de la position couchée (patte
horizontale).

\begin{figure}[H]
    \centering
    \resizebox{0.7\textwidth}{!}{%
    \begin{tikzpicture}[>=Stealth]
        % --- Bracelet gauche (vers la boucle) ---
        \fill[ifblue!35, rounded corners=2pt]
            (-3.0,0.65) rectangle (0,1.35);
        % Boucle de fermeture
        \fill[ifblue!50, rounded corners=2pt]
            (-3.6,0.5) rectangle (-3.0,1.5);
        \draw[ifblue!70, thick, rounded corners=1pt]
            (-3.45,0.7) rectangle (-3.15,1.3);

        % --- Bracelet droit (trous) ---
        \fill[ifblue!35, rounded corners=2pt]
            (4.0,0.65) rectangle (7.0,1.35);
        % Trous d'ajustement (ovales)
        \foreach \x in {4.6, 5.2, 5.8, 6.4} {
            \fill[white] (\x,1.0) ellipse (0.12 and 0.15);
        }

        % --- Boîtier central (vue de face) ---
        \fill[ifblue!50, rounded corners=4pt]
            (0,0) rectangle (4.0,2.0);
        % Bord supérieur légèrement plus clair
        \fill[ifblue!38, rounded corners=3pt]
            (0.1,1.55) rectangle (3.9,1.95);

        % Étiquette blanche
        \fill[white, rounded corners=2pt]
            (0.6,0.35) rectangle (3.4,1.4);
        \node[font=\small\bfseries, black] at (2.0,1.05) {IceQube};
        \node[font=\scriptsize, black!60] at (2.0,0.7) {IceRobotics};

        % --- Annotations ---
        \draw[->, black!55, thick]
            (4.5,2.5) -- (3.5,2.05)
            node[pos=0, anchor=south west, font=\scriptsize, black!70]
            {Boîtier étanche};

        \draw[->, black!55, thick]
            (-3.6,0.0) -- (-3.45,0.45)
            node[pos=0, anchor=north, font=\scriptsize, black!70]
            {Boucle};

        \draw[->, black!55, thick]
            (7.0,0.0) -- (5.8,0.6)
            node[pos=0, anchor=north, font=\scriptsize, black!70]
            {Trous d'ajustement};
    \end{tikzpicture}%
    }%
    \caption{Schéma du capteur IceQube.}
    \label{fig:iceqube_capteur}
\end{figure}

\begin{figure}[H]
    \centering
    \resizebox{0.4\textwidth}{!}{%
    \begin{tikzpicture}[>=Stealth]
        % --- Patte (forme simple, allongée verticalement) ---
        \fill[brown!12, rounded corners=5pt]
            (0,0) -- (0,6.0) -- (0.4,6.5) -- (1.4,6.5) -- (1.8,6.0)
            -- (1.8,0) -- (1.5,-0.4) -- (0.3,-0.4) -- cycle;
        \draw[brown!30, thick, rounded corners=5pt]
            (0,0) -- (0,6.0) -- (0.4,6.5) -- (1.4,6.5) -- (1.8,6.0)
            -- (1.8,0) -- (1.5,-0.4) -- (0.3,-0.4) -- cycle;

        % Jarret (ligne courte, texte à droite)
        \draw[brown!35, thick] (-0.2,4.5) -- (2.2,4.5);
        \node[font=\scriptsize, brown!55, anchor=west] at (2.4,4.5) {Jarret};

        % Zone métatarse (lignes courtes à droite seulement)
        \draw[brown!25, dashed] (1.9,1.0) -- (2.8,1.0);
        \draw[brown!25, dashed] (1.9,3.2) -- (2.8,3.2);
        \draw[decorate, decoration={brace, amplitude=4pt, mirror}, brown!40]
            (2.9,1.0) -- (2.9,3.2)
            node[midway, right=5pt, font=\scriptsize, brown!55] {Métatarse};

        % --- Bracelet autour de la patte ---
        \fill[ifblue!35, rounded corners=2pt]
            (-0.6,1.7) rectangle (0,2.5);
        \fill[ifblue!35, rounded corners=2pt]
            (1.8,1.7) rectangle (2.4,2.5);
        % Boucle côté gauche
        \fill[ifblue!50, rounded corners=1pt]
            (-0.8,1.8) rectangle (-0.6,2.4);

        % --- Capteur (sur la face avant de la patte) ---
        \fill[ifblue!50, rounded corners=3pt]
            (0.0,1.6) rectangle (1.8,2.6);
        \fill[white, rounded corners=1pt]
            (0.25,1.75) rectangle (1.55,2.45);
        \node[font=\scriptsize\bfseries, black] at (0.9,2.1) {IceQube};

        % --- Axes de l'accéléromètre (en bas à gauche, bien isolés) ---
        \coordinate (AX) at (-2.2,0.0);

        % z vers le haut
        \draw[->, rawred, very thick] (AX) -- ++(0,1.4)
            node[above, font=\small, rawred] {$z$};

        % x vers la droite
        \draw[->, rulegreen, very thick] (AX) -- ++(1.2,0)
            node[right, font=\small, rulegreen] {$x$};

        % y vers l'avant (hors plan)
        \draw[->, ifblue!80!black, very thick] (AX) -- ++(-0.7,-0.7)
            node[below left, font=\small, ifblue!80!black] {$y$};

        % --- Gravité (à droite, au-dessus du jarret) ---
        \draw[->, black!45, thick, dashed] (3.8,6.0) -- ++(0,-1.3)
            node[midway, right=2pt, font=\scriptsize, black!50] {$\vec{g}$};
    \end{tikzpicture}%
    }%
    \caption{Positionnement du capteur IceQube.}
    \label{fig:iceqube_position}
\end{figure}

Le capteur se fixe au métatarse d'une patte arrière par un bracelet à boucle.
L'accéléromètre interne mesure l'accélération sur trois axes. La composante gravitationnelle
contribue à distinguer la position couchée de la position debout, tandis que les variations
d'accélération alimentent les indicateurs d'activité. L'algorithme propriétaire agrège les
mesures brutes; le présent projet exploite des intervalles de 15~minutes. Cette résolution
est un choix analytique du mémoire et non une fréquence cliniquement optimale démontrée.
Des études de validation de capteurs portés rapportent une forte concordance entre certaines
mesures de posture et l'observation directe. \citet{alsaaod2012electronic} rapportent
une forte concordance entre les mesures d'accéléromètres de patte et l'observation directe
pour la classification couché/debout, tandis que \cite{blackie2011lying} montrent que ces
mesures permettent de caractériser les différences de comportement couché associées à la
boiterie.

Chaque intervalle de 15~minutes génère sept colonnes de données, dont cinq sont utilisées dans la
chaîne de traitement:
\begin{itemize}
    \item \textit{Steps}: nombre de pas détectés par l'accéléromètre;
    \item \textit{Motion Index}: indice synthétique d'activité locomotrice (unité
    propriétaire IceRobotics, dérivée des variations d'accélération);
    \item \textit{Lying Time}: durée cumulée en position couchée (en minutes, bornée à 15);
    \item \textit{Standing Time}: durée cumulée en position debout (en minutes, bornée à 15);
    \item \textit{Transitions}: nombre total de changements de posture
    couché$\leftrightarrow$debout.
\end{itemize}
Le \textit{Motion Index} a d'ailleurs déjà permis de détecter des comportements précis au
moyen de seuils validés~: \citet{mckay2024play} repèrent ainsi le jeu chez le veau sevré
équipé d'un accéléromètre IceTag. Cela confirme la valeur discriminante du signal, tout en
rappelant qu'un même indice recouvre des comportements variés et n'est donc pas spécifique à
la boiterie.

Les colonnes supplémentaires (\textit{Transitions Up} et \textit{Transitions Down})
détaillent la direction des transitions mais ne sont pas exploitées directement par la
chaîne de traitement, qui utilise leur somme.

Les données agrégées sont stockées dans la mémoire interne du capteur, puis récupérées
périodiquement. Elles sont accessibles via la plateforme web CowAlert (IceRobotics), qui
offre une interface de visualisation en ligne, et peuvent être exportées au format tabulaire
(CSV), directement exploitable par la chaîne de traitement décrite dans ce mémoire.

Le modèle IceTag, prédécesseur de l'IceQube, fournit des variables similaires (pas, temps
couché, transitions) et a été utilisé dans plusieurs travaux sur le comportement et la
locomotion bovins \citep{beer2016objective,alsaaod2012electronic}. Le choix de l'IceQube
est d'abord déterminé par la disponibilité du corpus. Il est cohérent avec la place
importante des accéléromètres portés dans la littérature sur la surveillance locomotrice
\citep{oleary2020accelerometers,bewley2008review}.

\cite{ito2010lying} montrent que les vaches boiteuses modifient significativement leur
comportement de couchage (durée allongée, fréquence de transitions réduite), ce qui
justifie l'inclusion de ces variables dans la chaîne de traitement. \cite{nechanitzky2016analysis}
montrent également que des lésions de la corne sont associées à des changements de
comportement mesurables. Ces résultats justifient l'étude de signaux comportementaux,
sans établir à eux seuls un délai universel de détection préclinique.

% =====================================================================
\section{Hétérogénéité des protocoles~: un obstacle à la comparaison}
% =====================================================================

Les revues spécialisées sur la détection automatisée de boiterie convergent vers une
limite récurrente: l'hétérogénéité des protocoles. Le type de capteur, sa position sur
l'animal, la définition opérationnelle de la boiterie, la taille des cohortes et l'horizon
d'évaluation varient fortement d'une étude à l'autre, ce qui complique la comparaison
directe des performances publiées, comme le soulignent les travaux de
\citet{oleary2020accelerometers} et de \citet{alsaaod2019automatic}.

Les travaux représentatifs du Tableau~\ref{tab:revue_ciblee} montrent pourquoi les
performances numériques ne doivent pas être juxtaposées sans précaution: les définitions
cliniques, les populations, les modalités de capteurs, les découpages d'évaluation et les
unités statistiques diffèrent. La présente section examine donc les protocoles plutôt qu'un
classement de sensibilités.
La revue d'O'Leary et al. \citep{oleary2020accelerometers} sur les accéléromètres pour la
détection de boiterie confirme cette dispersion et souligne l'absence de standardisation
dans les protocoles d'évaluation.
Alsaaod et al. \citep{alsaaod2019automatic} soulignent en particulier l'absence
de véritables études de longue durée ayant abouti à un système d'alerte précoce valide de
bout en bout.

Une étude longitudinale menée dans six fermes allemandes combine activité, couchage,
production laitière et caractéristiques individuelles dans un modèle à effets mixtes
\citep{lavrova2023lameness}. Elle illustre l'intérêt d'informations hétérogènes, mais son
protocole supervisé et ses données contextuelles diffèrent de ceux du présent mémoire.

Enfin, la qualité de la mesure amont reste déterminante: l'échantillonnage, l'édition des
épisodes et la représentation des comportements influencent directement la fiabilité des
variables dérivées \citep{ledgerwood2010lying,balasso2023behaviour}.

Concrètement, l'hétérogénéité des protocoles se manifeste sur au moins cinq dimensions:
\begin{enumerate}
    \item \textit{Définition du cas}: classe clinique, score locomoteur seuil ou probabilité
    de boiterie selon les études, d'où des phénomènes mesurés et des sensibilités
    difficilement comparables;
    \item \textit{Instrumentation}: position du capteur (patte
    \citep{alsaaod2012electronic,thorup2015lameness}, collier
    \citep{vanhertem2013lameness}, vision dorsale \citep{viazzi2013analysis}), fréquence
    d'échantillonnage et signaux
    contextuels varient, produisant des variables de nature différente;
    \item \textit{Granularité analytique}: prédiction instantanée, journalière, par épisode
    ou par vache, avec des exigences de précision temporelle très différentes;
    \item \textit{Objectif métier}: classification clinique, alerte précoce, filtrage ou
    score de risque, chaque objectif imposant des coûts d'erreur différents;
    \item \textit{Évaluation}: sensibilité, spécificité, AUC, taux d'alertes, IoU ou
    robustesse inter-fermes, un même système paraissant excellent ou médiocre selon la
    métrique.
\end{enumerate}

Autrement dit, deux articles peuvent rapporter une ``bonne performance'' tout en répondant à
des problèmes différents. Cette observation justifie le choix de consacrer un chapitre entier
aux protocoles de validation et à leurs artefacts, plutôt que de se limiter à un score unique
agrégé.

\FloatBarrier

% =====================================================================
\section{Détection temporelle de changements persistants}

La boiterie n'est pas nécessairement représentée par un point extrême isolé. Les changements
rapportés dans les signaux portés concernent souvent une tendance, une rupture ou une
modification persistante par rapport au comportement habituel de l'animal
\citep{alsaaod2019automatic,paudyal2021lying}. Cette structure temporelle favorise les
méthodes capables d'accumuler une évidence unilatérale plutôt que de classer indépendamment
chaque intervalle.

Les sommes cumulées proposées par \citet{page1954continuous} constituent une famille
classique de détection de changement. Appliquées à une combinaison de déficits normalisés,
elles permettent d'oublier les fluctuations brèves, d'accumuler une baisse soutenue et de
déclencher une alerte lorsque l'évidence franchit un seuil. Dans un contexte bovin, cette
logique doit être associée à une période de référence individuelle, au cycle journalier, à la qualité des
données et à la concordance de plusieurs familles de signaux. Elle reste toutefois une
détection comportementale: seule une validation avec examens locomoteurs et diagnostics
contemporains peut établir la relation clinique.

% =====================================================================
\section{Synthèse critique}
% =====================================================================

La littérature et les guides pratiques indiquent trois exigences pour une solution de
surveillance associée au risque de boiterie:
\begin{enumerate}
  \item \textit{Réactivité opérationnelle}: repérer des changements persistants assez tôt
  pour prioriser l'observation;
  \item \textit{Maîtrise de la charge d'alerte}: limiter les notifications sans suite opérationnelle
  et la fatigue d'alerte, qui dégradent rapidement la confiance des utilisateurs
  \citep{rutten2013invited};
  \item \textit{Interprétabilité}: fournir des raisons explicites utiles à la vérification pour
  chaque alerte, afin de faciliter la prise de décision terrain.
\end{enumerate}

En pratique, les méthodes purement pilotées par les données peuvent manquer d'explicabilité pour
l'opérationnel, tandis que les approches purement à règles peuvent manquer de flexibilité
face à la variabilité inter-animaux. Un compromis robuste consiste à coupler détection
d'anomalies et contraintes métier explicites. Ce couplage cumule la
capacité de généralisation des algorithmes non supervisés et des gardes-fous
interprétables. La littérature récente confirme cette orientation: les systèmes
d'élevage de précision les plus prometteurs combinent des signaux hétérogènes avec des
logiques de décision interprétables, plutôt que de s'appuyer sur un modèle unique en boîte
noire \citep{palma2025scoping,michelena2025plf_review}.

L'hétérogénéité des protocoles évoquée à la section précédente a une conséquence directe
sur l'évaluation des systèmes: de nombreux travaux comparent des capteurs, des fermes, des
labels et des définitions de cas différents, ce qui complique l'appréciation d'une chaîne
déployable de bout en bout. Dans ce contexte, il existe une place claire pour des
contributions moins théoriques mais plus systémiques: une chaîne reproductible, une
validation explicite sur plusieurs niveaux et une traçabilité des artefacts expérimentaux.

La revue de littérature conduit donc à un enseignement central: la difficulté n'est pas
seulement de proposer un détecteur performant, mais de construire une chaîne de décision
cohérente, reproductible et défendable du point de vue métier. Une telle exigence justifie
la structure du présent mémoire, qui traite l'algorithme, les règles métier, l'interface et
la validation comme un ensemble cohérent plutôt que comme des blocs séparés.

% =====================================================================
\section{Positionnement du mémoire}
% =====================================================================

Le présent mémoire retient d'abord une approche temporelle directionnelle orientée vers les
baisses d'activité, puis examine une extension bidirectionnelle. Cette hiérarchie reflète le
niveau de preuve: la branche HYPO est soumise à une validation attribuable, une ablation et
une analyse de sensibilité complètes; la branche d'instabilité teste une hypothèse
complémentaire qui reste exploratoire. La contribution relève de l'ingénierie scientifique:
construire une chaîne reproductible qui compare chaque vache à sa propre période de référence et évaluer
honnêtement ce que permettent des données d'activité non étiquetées. Concrètement, le
mémoire livre:
\begin{enumerate}
  \item une branche HYPO évaluée sur 44 événements postérieurs à la période de référence, avec exécution propre de
  référence, ablation contre IF et LOF et analyse OFAT de vingt variantes;
  \item une branche d'instabilité exploratoire et cinq stratégies de fusion, dont une règle
  hiérarchique séparant surveillance et notification de vérification;
  \item une campagne étendue de 132 événements incluant contrôles brefs et facteurs
  confondants;
  \item une évaluation du temps d'exécution de la chaîne complète sur les 28 vaches;
  \item une analyse complémentaire IceTag--SLS synchronisée, explicitement interprétée comme
  exploratoire et non concluante;
  \item une traçabilité par tests et artefacts tabulaires reproductibles.
\end{enumerate}

La portée reste clairement délimitée: la sortie est une alerte de dégradation
comportementale à vérifier, et non une classification ou un diagnostic vétérinaire de
boiterie. Cette distinction est fondamentale et cohérente avec les recommandations de la
littérature \citep{alsaaod2019automatic}.

Le chapitre suivant détaille donc les données utilisées, les transformations appliquées aux
signaux capteurs, le choix des variables et le protocole expérimental retenu. L'objectif n'est
pas seulement de décrire une implémentation, mais de montrer comment la revue précédente se
traduit en décisions méthodologiques explicites.
